\documentclass[11pt]{amsart}
\usepackage{amsmath,amssymb,amsthm}
\usepackage[margin=1.1in]{geometry}
\usepackage{booktabs}
\usepackage[colorlinks=true,linkcolor=blue,citecolor=blue,urlcolor=blue]{hyperref}
\emergencystretch=2.5em
\raggedbottom

\newtheorem{theorem}{Theorem}[section]
\newtheorem{lemma}[theorem]{Lemma}
\newtheorem{proposition}[theorem]{Proposition}
\newtheorem{corollary}[theorem]{Corollary}
\theoremstyle{remark}
\newtheorem{remark}[theorem]{Remark}

\newcommand{\Z}{\mathbb{Z}}
\newcommand{\Q}{\mathbb{Q}}
\newcommand{\HH}{\mathbb{H}}
\newcommand{\e}{\varepsilon}
\newcommand{\dd}{\mathrm{d}}
\DeclareMathOperator{\gen}{gen}
\DeclareMathOperator{\vol}{vol}

\title[Theta series of the odd Heegner discriminants]{The theta series $\theta_D(z)\theta_D(pz)$\\ for the odd Heegner discriminants}
\author{Ruqing Chen}
\date{\today\ (chapter draft)}

\begin{document}
\begin{abstract}
For the seven odd discriminants of class number one, $D\in\{3,7,11,19,43,67,163\}$, the weight-one theta series $\theta_D$ of the norm form $N_D=x^2+xy+\frac{D+1}4y^2$ is the Eisenstein series $E_1(z,\chi_{-D})$, and $\Theta^{(D)}_{1,p}=\theta_D(z)\theta_D(pz)$ is a form of weight $2$ for $\Gamma_0(Dp)$, a square-free level whose four cusps are all Atkin--Lehner. We prove, uniformly in $D$: the Siegel main term with closed-form local densities at $D$ and at $p$; the signs $\Theta|W_p=\chi(p)\Theta$, $\Theta|W_D=-\chi(p)\Theta$, $\Theta|W_{Dp}=-\Theta$, hence the support of the cuspidal part on the newforms with $a_\varphi(D)=\chi_{-D}(p)$, $a_\varphi(p)=-\chi_{-D}(p)$ (no inert vanishing), on the oldforms from level $p$ with $a_f(p)=-\chi_{-D}(p)$, and on the oldforms from level $D$ only when $\chi_{-D}(p)=+1$; the Rankin--Selberg formula $\langle\Theta,\varphi\rangle=C_D\chi(p)\eta_p(\varphi)L(1,\varphi)L(1,\varphi\otimes\chi_{-D})$ with $C_D=w_D^2\sqrt D/(8\pi^2)$; and the growth $\|G^{(D)}_{1,p}\|^2\sim\frac6{\pi(D+1)}(\log p)^2$, with the second term of the two wide cusps. The constants are those of the analytic class number formula with $h=1$: $w_D^2L(1,\chi_{-D})^2=4\pi^2/D$. Everything is verified numerically for $D=3,7,11$ ($43$ newforms, Rankin--Selberg ratios within $0.5\%$, densities to $10^{-16}$).
\end{abstract}
\maketitle

\section{Setting}
Let $D$ be one of the seven odd Heegner discriminants, $K=\Q(\sqrt{-D})$, $\chi=\chi_{-D}=\bigl(\frac{-D}\cdot\bigr)$ the Kronecker symbol, $w_D$ the number of units ($w_3=6$, $w_D=2$ otherwise), $N_D(x,y)=x^2+xy+\frac{D+1}4y^2$ the norm form (Gram matrix $\left(\begin{smallmatrix}2&1\\1&(D+1)/2\end{smallmatrix}\right)$ of determinant $D$), and
\begin{gather*}
r_D(n)=\#\{(x,y):\ N_D(x,y)=n\}=w_D\sum_{d\mid n}\chi(d)\quad(n\ge1),\\
\theta_D(z)=\sum_{n\ge0}r_D(n)e(nz)\in M_1(\Gamma_0(D),\chi),
\end{gather*}
the formula for $r_D$ being the class number one property. The Eisenstein series $E(z,s)=\sum_{\Gamma_\infty\backslash\Gamma_0(D)}\chi(d_\gamma)j(\gamma,z)^{-1}\operatorname{Im}(\gamma z)^s$ has, at $s=0$, the expansion $1+\frac2{L(0,\chi)}\sum\sigma_\chi(n)e(nz)$ with $L(0,\chi)=2h/w_D=2/w_D$ (Hecke), so $E(z,0)=\theta_D(z)$ without any dimension argument. For a prime $p\ne D$ put
\begin{gather*}
\Theta_{1,p}=\theta_D(z)\theta_D(pz)=E+G\in M_2(\Gamma_0(Dp)),\\
T(k)=\sum_{\alpha+p\beta=k}r_D(\alpha)r_D(\beta),\qquad a_G(k)=T(k)-M(k),
\end{gather*}
$E$ the Eisenstein and $G=G^{(D)}_{1,p}$ the cuspidal part. The level $Dp$ is square-free; $\Gamma_0(Dp)$ has the four cusps $W_Q\infty$, $Q\in\{1,D,p,Dp\}$, of widths $Dp,p,D,1$, and no cusp of the type ``denominator $2\delta$'' of Part~II. (The even discriminants $D=4,8$, with $\theta_4=\theta^2$ and $\theta_8=\theta(z)\theta(2z)$, have levels $4p$, $8p$ with non-Atkin--Lehner cusps; $D=4$ is Parts~I--III, and $D=8$ is not treated here.) By $|_k$ we denote the weight-$k$ slash operator normalised by $(\det\gamma)^{k/2}$, so that scalars act trivially and $f(mz)=m^{-k/2}(f|_kB_m)(z)$, $B_m=\left(\begin{smallmatrix}m&0\\0&1\end{smallmatrix}\right)$.

\section{The Siegel main term}
\begin{theorem}\label{thm:siegel}
$M(k)=\dfrac{4\pi^2k}{Dp}\prod_\ell\beta_\ell(k)$, where $\beta_\ell(k)=(1-\ell^{-2})\sigma_{-1}(\ell^{v_\ell(k)})$ for every prime $\ell\nmid Dp$, including $\ell=2$, and, with $v=v_\ell(k)$,
\[
\beta_p(k)=(1+\chi(p))-\chi(p)\Bigl(1+\frac1p\Bigr)p^{-v},\qquad \beta_D(k)=(1+\e)-\e\Bigl(1+\frac1D\Bigr)D^{-v},
\]
where $\e=\bigl(\frac{-p}D\bigr)=-\chi_{-D}(p)$.
\end{theorem}
\begin{proof}
The quaternary form $Q=N_D(x,y)+pN_D(z,w)$ has Gram determinant $D^2p^2$, and Siegel's theorem gives $r(\gen Q,k)=4\pi^2k\prod_\ell\beta_\ell(k)/\sqrt{D^2p^2}$ with $\beta_\ell(k)=\lim_a\ell^{-3a}\#\{x\bmod\ell^a:\ Q(x)\equiv k\}$; the genus theta series is the Eisenstein part of $\Theta_{1,p}$. For $\ell\nmid Dp$ the form is unimodular over $\Z_\ell$ with square discriminant, which gives the standard density; this includes $\ell=2$ because $N_D$ is even unimodular over $\Z_2$, $pN_D\cong N_D$ there ($p$ a unit, the Arf invariant being insensitive to unit scaling), and $N_D\oplus N_D$ has Arf invariant $0$, i.e.\ is the sum of two hyperbolic planes, whatever the class of $N_D$ modulo $2$ ($D\equiv3$ or $7\pmod8$).

At $\ell=p$: the solutions with $(x,y)\not\equiv(0,0)\pmod p$ are non-singular in $(x,y)$ (the gradient of $N_D$ does not vanish there, $N_D$ being non-degenerate modulo $p$), so each solution modulo $p$ lifts to $p^{a-1}$ solutions modulo $p^a$; modulo $p$ the equation $N_D(x,y)\equiv t$ has $p-\chi(p)$ solutions for $p\nmid t$ (the number of elements of norm $t$ in $\mathbb F_{p^2}$ or $\mathbb F_p\times\mathbb F_p$) and $1+(p-1)(1+\chi(p))$ for $p\mid t$. With $(z,w)$ free, this gives $\beta_p(k)=1-\chi(p)/p$ for $p\nmid k$ and the contribution $(1+\chi(p))(1-1/p)$ for $p\mid k$. The solutions with $x\equiv y\equiv0\pmod p$, $x=px'$, $y=py'$, satisfy $pN_D(x',y')+N_D(z,w)\equiv k/p\pmod{p^{a-1}}$, the same form with the blocks interchanged; counting $(x',y')$ modulo $p^{a-1}$ and the $p^2$ lifts of $(z,w)$ gives the contribution $\beta_p(k/p)/p$. Hence $\beta_p(k)=(1+\chi(p))(1-1/p)+\beta_p(k/p)/p$ for $p\mid k$, which iterates to the closed form.

At $\ell=D$: since $2$ is invertible, $N_D(x,y)=(x+y/2)^2+\frac D4y^2\cong u^2+Dv^2$ over $\Z_D$, so $Q\cong u^2+ps^2+D(v^2+pt^2)$. The solutions with $(u,s)\not\equiv0\pmod D$ are non-singular in $(u,s)$; modulo $D$ the binary form $u^2+ps^2$, of discriminant $-4p$, represents $t\not\equiv0$ exactly $D-\e$ times and $t\equiv0$ exactly $1+(D-1)(1+\e)$ times, $\e=(-p/D)$. The solutions with $u\equiv s\equiv0$ reduce to $D(u'^2+ps'^2)+(v^2+pt^2)\equiv k/D\pmod{D^{a-1}}$, the same form, with the factor $D^2/D^3=1/D$. Hence $\beta_D(k)=1-\e/D$ for $D\nmid k$ and $\beta_D(k)=(1+\e)(1-1/D)+\beta_D(k/D)/D$ for $D\mid k$. Finally $(-p/D)=(-1/D)(p/D)=-(p/D)=-(-D/p)=-\chi_{-D}(p)$ by quadratic reciprocity, $D\equiv3\pmod4$.
\end{proof}

For $D=3$ this is Proposition~2.1 of the working notes; $1+1/D=4/3$ there. The script \texttt{heegner\_theta\_unified.py} checks the closed forms against direct counts modulo $D^a$ and $p^a$ ($D=3,7,11$, $p\le13$; relative error $2\cdot10^{-16}$), and the whole main term against the known cusp forms: at the level $15$, $T-M=\frac92f_{15}$; at the level $14$, $T-M=\frac23f_{14}$ ($f_{14}=\eta(z)\eta(2z)\eta(7z)\eta(14z)$); at the level $22$, $T-M=0$ identically (Section~\ref{sec:support} explains why); at the levels $21,33,35,39,55,77,91,143$ the mean square of $T-M$ grows like $Y^2$, as for a cusp form.

\section{Atkin--Lehner signs and the spectral support}\label{sec:support}
\begin{theorem}\label{thm:AL}
$\Theta_{1,p}|_2W_p=\chi(p)\Theta_{1,p}$, $\Theta_{1,p}|_2W_D=-\chi(p)\Theta_{1,p}$ and $\Theta_{1,p}|_2W_{Dp}=-\Theta_{1,p}$, for any choice of the Atkin--Lehner matrices of level $Dp$.
\end{theorem}
\begin{proof}
Write $W_p=\left(\begin{smallmatrix}pa&b\\Dpc&pd\end{smallmatrix}\right)$, $pad-Dbc=1$. Then $W_p=\gamma B_p$ with $\gamma=\left(\begin{smallmatrix}a&b\\Dc&pd\end{smallmatrix}\right)\in\Gamma_0(D)$, so $\theta_D|_1W_p=\chi(pd)\,\theta_D|_1B_p=\chi(pd)p^{1/2}\theta_D(pz)$; and $B_pW_p=p\left(\begin{smallmatrix}pa&b\\Dc&d\end{smallmatrix}\right)$ with the second factor in $\Gamma_0(D)$, so $\theta_D(p\,\cdot)|_1W_p=p^{-1/2}\chi(d)\theta_D$. The product is $\chi(p)\chi(d)^2\Theta_{1,p}=\chi(p)\Theta_{1,p}$.

For $W_D$ we use the transformation of $\theta_D$ under its own Atkin--Lehner matrices. The dual of the lattice $(\mathcal O_K,N_D)$ is $\mathcal O_K/\sqrt{-D}$, so Poisson summation gives $\theta_D(-1/z)=D^{-1/2}(-iz)\theta_D(z/D)$ and $\theta_D(-1/(Dz))=-i\sqrt Dz\,\theta_D(z)$, i.e.\ $\theta_D|_1W_0=-i\theta_D$ for $W_0=\left(\begin{smallmatrix}0&-1\\D&0\end{smallmatrix}\right)$. Every matrix $W=\left(\begin{smallmatrix}Da'&b'\\Dc'&Dd'\end{smallmatrix}\right)$ of determinant $D$ is $W_0\gamma'$ with $\gamma'=\left(\begin{smallmatrix}c'&d'\\-Da'&-b'\end{smallmatrix}\right)\in\Gamma_0(D)$, whence $\theta_D|_1W=-i\chi(-b')\theta_D$. Now $W_D=\left(\begin{smallmatrix}Da&b\\Dpc&Dd\end{smallmatrix}\right)$, $Dad-pbc=1$, is such a matrix with $b'=b$, and $B_pW_D=\left(\begin{smallmatrix}Da&pb\\Dc&Dd\end{smallmatrix}\right)B_p$ exhibits $\theta_D(p\,\cdot)|_1W_D=p^{-1/2}(\theta_D|_1W')|_1B_p$ with $W'$ such a matrix with $b'=pb$. Therefore $\Theta_{1,p}|_2W_D=(-i)^2\chi(-b)\chi(-pb)\Theta_{1,p}=-\chi(p)\Theta_{1,p}$, using $(b,D)=1$. Finally $W_{Dp}\equiv W_DW_p$ modulo $\Gamma_0(Dp)$ gives $-\chi(p)\cdot\chi(p)=-1$. The signs are checked numerically at $D=3,7,11$, $p\le13$, to $10^{-15}$.
\end{proof}

\begin{theorem}[Spectral support]\label{thm:support}
Let $\varphi$ be a newform of weight $2$ and level $Dp$. If $\langle\Theta_{1,p},\varphi\rangle\ne0$ then
\[
\eta_p(\varphi)=\chi(p),\quad\eta_D(\varphi)=-\chi(p),\qquad\text{i.e.}\qquad a_\varphi(D)=\chi_{-D}(p),\quad a_\varphi(p)=-\chi_{-D}(p),
\]
and the root number of $\varphi$ is $+1$. For a newform $f$ of level $p$ and $d\mid D$, $\langle\Theta_{1,p},f(d\,\cdot)\rangle\ne0$ requires $\eta_p(f)=\chi(p)$, i.e.\ $a_f(p)=-\chi_{-D}(p)$ (root number $+1$). For a newform $f$ of level $D$ and $d\mid p$, $\langle\Theta_{1,p},f(d\,\cdot)\rangle=C_D\bigl(1+\chi(p)\bigr)L(1,f)L(1,f\otimes\chi)/d$ vanishes when $\chi_{-D}(p)=-1$.
\end{theorem}
\begin{proof}
$\langle\Theta,\varphi\rangle=\langle\Theta|W,\varphi|W\rangle$ for $W=W_p,W_D$ and Theorem~\ref{thm:AL} give $\chi(p)\eta_p(\varphi)=1$ and $-\chi(p)\eta_D(\varphi)=1$; $a_\varphi(\ell)=-\eta_\ell(\varphi)$ for $\ell\,\|\,Dp$, and the root number is $-\eta_D\eta_p$. For $f$ of level $p$, Lemma~6.6 of Part~II (with $Q=p$, $Q_1=p$, $Q_2=Q_3=1$) gives $f(d\,\cdot)|_2W_p=\eta_p(f)f(d\,\cdot)$, and the same argument applies; $W_D$ exchanges $f(z)$ and $f(Dz)$ up to a factor and gives no further condition. For $f$ of level $D$ the formula is the oldform computation in the proof of Part~II, Theorem~6.7(ii), with $D$ in the role of $4$ and $p$ in the role of $n$ (the level of $f$ is ``$2^e$'' with $M_1=1$, and $d\mid p$ is the ``$d_{\mathrm o}$''): the second bracket there is $\sum_{Q_2\mid d}\chi(Q_2)$ and the third is $\prod_{\ell\mid p/d}(1+\chi(\ell))$, so the product is $1+\chi(p)$ in both cases $d=1$, $d=p$.
\end{proof}

In contrast with $D=4$, where every newform of level $4n$ has $\eta_4=-1$ (Part~II, Proposition~6.17), so that the newform support is empty when $\chi_{-4}(n)=-1$, here both eigenvalues of $W_D$ occur among the newforms of level $Dp$, since the local component at $D$ may be the Steinberg representation or its unramified quadratic twist: there is no inert vanishing, and for every $p$ the newform support is the subspace $\{a_\varphi(D)=\chi(p),\ a_\varphi(p)=-\chi(p)\}$, cut down further by $L(1,\varphi)L(1,\varphi\otimes\chi)\ne0$. This subspace is not exactly half of the new subspace: its dimension is $\frac14(\dim S^{\mathrm{new}}+\chi(p)\operatorname{tr}W_p-\chi(p)\operatorname{tr}W_D-\operatorname{tr}W_{Dp})$, the traces taken on $S_2^{\mathrm{new}}(\Gamma_0(Dp))$, and the traces of the Atkin--Lehner involutions are sums of class numbers of size $O(\sqrt{Dp}\log Dp)$ (Yamauchi's trace formula), so it is asymptotically one quarter of the new subspace, i.e.\ half of the forms with root number $+1$. [Citation to be checked.] The data: for $D=3$, $p\in\{5,7,11,13,17,19,23,29,31,37\}$ the numbers of newform components of $G_{1,p}$ are $1,1,1,2,2,2,2,3,3,4$ against $\dim S^{\mathrm{new}}=1,1,1,1,3,3,3,5,5,7$; for $D=7$, $p\in\{5,11,13,17,19\}$: $2,3,3,5,4$; for $D=11$, $p\in\{5,7,13,17,19\}$: $2,3,6,6,9$. At the level $77$ the same three newforms carry both $G^{(7)}_{1,11}$ and $G^{(11)}_{1,7}$ ($a(7)=+1=\chi_{-7}(11)=-\chi_{-11}(7)$, $a(11)=-1$), with the shares $0.001,0.019,0.981$ and $0.502,0.495,0.005$ respectively, in accordance with the two different twists in Theorem~\ref{thm:RS}. The oldforms from level $D$ exist for $D\in\{11,19,43,67,163\}$ (genus $1,1,3,5,13$) and not for $D=3,7$; at the level $22$ ($\chi_{-11}(2)=-1$) the cuspidal part vanishes identically, $\theta_{11}(z)\theta_{11}(2z)=E$, while at the level $55$ ($\chi_{-11}(5)=+1$) the component of $f_{11}$ is $f_{11}(z)+5f_{11}(5z)$ up to a scalar and carries $25\%$ of the norm.

\section{The Rankin--Selberg formula}
\begin{theorem}\label{thm:RS}
For a newform $\varphi$ of weight $2$ and level $Dp$,
\[
\langle\Theta_{1,p},\varphi\rangle=C_D\,\chi(p)\,\eta_p(\varphi)\,L(1,\varphi)\,L(1,\varphi\otimes\chi_{-D}),\qquad C_D=\frac{w_D}{4\pi L(1,\chi_{-D})}=\frac{w_D^2\sqrt D}{8\pi^2}.
\]
Thus $C_3=\frac{9\sqrt3}{2\pi^2}=0.7897$, $C_D=\frac{\sqrt D}{2\pi^2}$ for $D\ge7$ ($C_7=0.1340$, $C_{11}=0.1680$), against $C_4=\frac4{\pi^2}$.
\end{theorem}
\begin{proof}
The proof of Part~II, Theorem~6.7(i) carries over: $\theta_D=E(z,0)$; the cusps of $\Gamma_0(Dp)$ over $\infty$ of $\Gamma_0(D)$ are $W_Q\infty$, $Q\mid p$; unfolding gives
\[
I(s)=\sum_{Q\mid p}\chi(Q)\,\Gamma(s+1)(4\pi)^{-s-1}\sum_{j\ge1}\frac{r_D(j)\,c_Q((p/Q)j)}{((p/Q)j)^{s+1}},
\]
with $\chi(a)\chi(d)=\chi(Q)$ from $Qad\equiv1\pmod D$, and $c_Q(k)=\eta_Q(\varphi)a_\varphi(k)$. The Dirichlet series is $w_D\,L(w,\varphi)L(w,\varphi\otimes\chi)/\bigl(L(2w-1,\chi)\prod_{\ell\mid p}(1-\chi(\ell)\ell^{1-2w})\bigr)$, the local factor at $\ell=D$ being $(1-a_\varphi(D)x)^{-1}=L_D(w,\varphi)L_D(w,\varphi\otimes\chi)$ because $\chi(D)=0$ and $\sigma_\chi(D^j)=1$. At $s=0$ the constant is $w_D/(4\pi L(1,\chi))$, and $\sum_{Q\mid p}\chi(Q)\eta_Q(\varphi)a_\varphi(p/Q)/(p/Q)=\chi(p)\eta_p(\varphi)(1-\chi(p)/p)$ cancels the local factor. The class number formula $L(1,\chi_{-D})=2\pi h/(w_D\sqrt D)$ with $h=1$ gives the second expression. Numerically, for the $43$ newform components of $G^{(D)}_{1,p}$ with $D=3,7,11$ and $p\le19$, the ratio of the two sides lies in $[0.9952,1.0048]$ (the twist $\varphi\otimes\chi_{-D}$ has level $D^2p$ and root number $+1$ in every case).
\end{proof}

\section{The norm}
\begin{theorem}\label{thm:norm}
Let $D_D(s)=\sum_kr_D(k)^2k^{-s}$. Then $D_D(s)=w_D^2\,\zeta(s)^2L(s,\chi)^2/\bigl((1+D^{-s})\zeta(2s)\bigr)$, and
\begin{gather*}
S_D(Y):=\sum_{k\ge1}\frac{r_D(k)^2}k e^{-k/Y}=\frac{12}{D+1}(\log Y)^2+c_1^{(D)}\log Y+c_2^{(D)}+O(Y^{-1/2}),\\
c_1^{(D)}=\frac{24}{D+1}\Bigl(\gamma+2\frac{L'}L(1,\chi)-2\frac{\zeta'}\zeta(2)+\frac{\log D}{D+1}\Bigr),
\end{gather*}
with $\frac{L'}L(1,\chi_{-D})=\gamma+\log2+\log\pi-\frac{w_D}2\sum_{a=1}^{D-1}\chi(a)\log\Gamma(a/D)$. Consequently the regions above height $1$ at the two wide cusps give
\begin{gather*}
\|G^{(D)}_{1,p}\|^2\ \ge\ \frac1{4\pi}\Bigl(S_D\Bigl(\frac{Dp}{4\pi}\Bigr)+S_D\Bigl(\frac p{4\pi}\Bigr)\Bigr)-O(\log p)=\frac6{\pi(D+1)}(\log p)^2+b_D\log p+O(1),\\
b_D=\frac{6}{\pi(D+1)}\log\frac D{16\pi^2}+\frac{c_1^{(D)}}{2\pi}\qquad(b_3=0.713,\ b_7=0.207,\ b_{11}=0.133),
\end{gather*}
and $\|G^{(D)}_{1,p}\|^2\sim\frac6{\pi(D+1)}(\log p)^2$ provided the upper bound of Part~II, Theorem~6.14, carries over (four Atkin--Lehner classes, no odd class).
\end{theorem}
\begin{proof}
$r_D/w_D=\sigma_\chi$ is multiplicative with $\sigma_\chi(\ell^j)=j+1$, $\frac12(1+(-1)^j)$, $1$ for $\chi(\ell)=1,-1,0$, so the Euler factor of $\sum\sigma_\chi(k)^2k^{-s}$ at $\ell\ne D$ is that of $\zeta^2L^2/\zeta(2s)$, and at $\ell=D$ it is $(1-x)^{-1}$ against $(1+x)/(1-x)$ for $\zeta^2L^2/\zeta(2s)$, whence the factor $(1+D^{-s})^{-1}$. At $s=1$: $w_D^2L(1,\chi)^2\cdot\frac{D}{D+1}\cdot\frac6{\pi^2}=\frac{4\pi^2}D\cdot\frac D{D+1}\cdot\frac6{\pi^2}=\frac{24}{D+1}$ by the class number formula, so $D_D(1+w)=\frac{24}{D+1}w^{-2}(1+a_1w+\dots)$ with $a_1=2\gamma+2\frac{L'}L(1,\chi)-2\frac{\zeta'}\zeta(2)+\frac{\log D}{D+1}$, and the residue computation of Part~II, Lemma~6.10 gives $S_D$. The value of $\frac{L'}L(1,\chi)$ follows from the functional equation $\Lambda(s)=(D/\pi)^{s/2}\Gamma(\frac{s+1}2)L(s,\chi)=\Lambda(1-s)$ (root number $+1$ for real odd $\chi$) and $L(s,\chi)=D^{-s}\sum_a\chi(a)\zeta(s,a/D)$ with $\zeta'(0,a)=\log\Gamma(a)-\frac12\log2\pi$: $\frac{L'}L(1)=-\log\frac D\pi-\frac12\psi(\tfrac12)-\frac12\psi(1)-\frac{L'}L(0)$ and $\frac{L'}L(0)=-\log D+\frac{w_D}2\sum_a\chi(a)\log\Gamma(a/D)$, since $L(0,\chi)=2/w_D$. The lower bound for the norm is Part~II, Lemma~6.9 at the cusps $W_1\infty$ and $W_D\infty$ (widths $Dp$ and $p$), which carry $\pm G_{1,p}$ by Theorem~\ref{thm:AL}, with the correction by the main term and by the coefficients $k\ge p$ bounded as in Part~II, Theorem~6.5(iii), and $\log\frac{Dp}{4\pi}+\log\frac p{4\pi}=2\log p+\log\frac D{16\pi^2}$.
\end{proof}

The growth constant is $\frac6{\pi(D+1)}$, decreasing in $D$ ($0.4775$, $0.2387$, $0.1592$ for $D=3,7,11$, against $\frac1\pi=0.3183$ for $D=4$, where the factor at the ramified prime is $\frac23$ instead of $\frac D{D+1}$); the value $\frac{\sqrt D}{2\pi}$ proposed for it is incorrect. What the class number formula fixes is the product $w_D^2L(1,\chi_{-D})^2=4\pi^2/D$, which makes the residue of $D_D$ independent of $w_D$; the remaining factor $D/(D+1)$ is the Euler factor at the ramified prime, where $\sigma_\chi(D^j)=1$.

\section{Numerical results}
\begin{center}\footnotesize\setlength{\tabcolsep}{4pt}
\begin{tabular}{c|cccc|ccc|cc}
\toprule
$D$ & $w_D$ & $L(1,\chi)$ & $\frac{L'}L(1,\chi)$ & $C_D$ & $\frac{12}{D+1}$ & $c_1^{(D)}$ & $c_2^{(D)}$ & $\frac6{\pi(D+1)}$ & fit $a$\\
\midrule
$3$ & $6$ & $0.604600$ & $0.368282$ & $0.789720$ & $3$ & $16.370$ & $13.059$ & $0.4775$ & $0.4815$\\
$7$ & $2$ & $1.187410$ & $0.015636$ & $0.134035$ & $1.5$ & $5.975$ & $1.450$ & $0.2387$ & $0.2408$\\
$11$ & $2$ & $0.947226$ & $-0.084218$ & $0.168022$ & $1$ & $3.497$ & $0.492$ & $0.1592$ & $0.1601$\\
\bottomrule
\end{tabular}
\end{center}
($L(1,\chi)$ and $\frac{L'}L(1,\chi)$ agree with mpmath to all digits shown; $c_2^{(D)}$ from $S_D(Y)$ at $Y=10^4,10^5$; ``fit $a$'' is the coefficient of $(\log p)^2$ in a three-parameter fit of $\Lambda_1(G_{1,p})$ over $p\in\{101,1009,10007,100003\}$.) The sandwich $\Lambda_1\le\|G\|^2\le\Lambda_{\sqrt3/2}$ over the four cusps, and the Rankin--Selberg estimate where it has converged ($p\le1009$):
\begin{center}\footnotesize\setlength{\tabcolsep}{4pt}
\begin{tabular}{r|rrr|rrr|rrr}
\toprule
 & \multicolumn{3}{c|}{$D=3$} & \multicolumn{3}{c|}{$D=7$} & \multicolumn{3}{c}{$D=11$}\\
$p$ & RS & $\Lambda_1$ & $\Lambda_{\sqrt3/2}$ & RS & $\Lambda_1$ & $\Lambda_{\sqrt3/2}$ & RS & $\Lambda_1$ & $\Lambda_{\sqrt3/2}$\\
\midrule
$101$ & $11.893$ & $11.621$ & $12.513$ & $5.437$ & $5.294$ & $5.774$ & $4.092$ & $3.969$ & $4.363$\\
$1009$ & $26.452$ & $25.966$ & $27.300$ & $12.556$ & $12.263$ & $13.013$ & $9.040$ & $8.783$ & $9.402$\\
$10007$ & & $45.385$ & $47.151$ & & $21.768$ & $22.785$ & & $15.277$ & $16.117$\\
$100003$ & & $69.921$ & $72.120$ & & $33.833$ & $35.116$ & & $23.479$ & $24.540$\\
\bottomrule
\end{tabular}
\end{center}
The exact norm by the scalloped Parseval formula of the working notes applies verbatim (four Atkin--Lehner classes) and has not been run for $D\ne4$.

\section{The remaining odd discriminants, and the even ones}
\subsection{$D=19,43,67,163$ (\texttt{results/heegner\_theta\_unified\_}$D$\texttt{.txt})}
The script was run for each of the four remaining discriminants (one at a time; the mpmath check of $\frac{L'}L(1,\chi)$, verified at $D\le19$, is skipped beyond). In all four cases the densities of Theorem~\ref{thm:siegel} agree with the counts modulo $D^2$ and $p^a$ to $2\cdot10^{-16}$, and the three Atkin--Lehner signs of Theorem~\ref{thm:AL} hold to $10^{-15}$ for $p\in\{2,5,7,11,13\}$. For $D=19$ ($N_{19}=x^2+xy+5y^2$, $L(1,\chi)=0.720731$, $\frac{L'}L(1,\chi)=-0.084914$, $C_{19}=0.220824$, $\frac{12}{D+1}=0.6$, $c_1^{(19)}=2.033$, $\frac6{\pi(D+1)}=0.0955$): the densities agree with the counts to $2\cdot10^{-16}$, the three Atkin--Lehner signs hold for $p\le13$, every newform component of $G_{1,p}$ for $p\in\{5,7,11,13,17\}$ ($23$ newforms) has $a(19)=\chi_{-19}(p)$, $a(p)=-\chi_{-19}(p)$, the Rankin--Selberg ratios lie in $[0.9965,1.0047]$, and the norm scan gives
\begin{center}\footnotesize
\begin{tabular}{r|rrr|rr}
\toprule
$p$ & RS & $\Lambda_1$ & $\Lambda_{\sqrt3/2}$ & $\frac6{20\pi}(\log p)^2$ & $\Lambda_1/(\log p)^2$\\
\midrule
$101$ & $3.402$ & $3.293$ & $3.637$ & $2.034$ & $0.155$\\
$1009$ & $6.533$ & $6.502$ & $7.023$ & $4.569$ & $0.136$\\
$10007$ & & $10.733$ & $11.429$ & $8.102$ & $0.127$\\
$100003$ & & $15.975$ & $16.843$ & $12.657$ & $0.121$\\
\bottomrule
\end{tabular}
\end{center}
with the fit $\Lambda_1=0.0961(\log p)^2+0.290\log p-0.09$ over $p\ge100$, against $0.0955$. For $D=43$ ($N_{43}=x^2+xy+11y^2$, $L(1,\chi)=0.479088$, $\frac{L'}L(1,\chi)=0.249482$, $C_{43}=0.332204$, $\frac{12}{D+1}=0.2727$, $c_1^{(43)}=1.255$): for $p=5,7,11$ (levels $215,301,473$, $\dim S_2=21,27,43$) the decomposition resolves $13$ newform components, all obeying the law, with Rankin--Selberg ratios in $[0.9936,1.0038]$ (at $p=5$ one numerically contaminated duplicate of a component, with $a(43)=-1.0001$, is mislabelled by the script; its clean copy obeys the law); for $p\ge13$ ($\dim S_2\ge49$) the Prony recurrence of order $\le11$ does not separate the span. For $D=67$ ($L(1,\chi)=0.383\ldots$, $\frac{L'}L=0.658364$, $c_1^{(67)}=1.093$) the components found at $p=5,7$ are mixtures and nothing is tested beyond the signs; for $D=163$ ($\frac{L'}L=2.168327$, $c_1^{(163)}=0.891$, $\dim S_2(\Gamma_0(815))=81$) the decomposition is out of reach of the method. The norm scans:
\begin{center}\footnotesize\setlength{\tabcolsep}{4pt}
\begin{tabular}{r|rrr|rrr|rrr}
\toprule
 & \multicolumn{3}{c|}{$D=43$, $\frac6{44\pi}=0.0434$} & \multicolumn{3}{c|}{$D=67$, $\frac6{68\pi}=0.0281$} & \multicolumn{3}{c}{$D=163$, $\frac6{164\pi}=0.0116$}\\
$p$ & RS & $\Lambda_1$ & $\Lambda_{\sqrt3/2}$ & RS & $\Lambda_1$ & $\Lambda_{\sqrt3/2}$ & RS & $\Lambda_1$ & $\Lambda_{\sqrt3/2}$\\
\midrule
$101$ & $3.241$ & $3.080$ & $3.342$ & $3.342$ & $3.005$ & $3.206$ & $3.487$ & $2.875$ & $2.976$\\
$1009$ & $4.873$ & $4.982$ & $5.339$ & $4.434$ & $4.531$ & $4.798$ & $3.969$ & $3.759$ & $3.885$\\
$10007$ & & $7.347$ & $7.796$ & & $6.304$ & $6.631$ & & $4.779$ & $4.924$\\
$100003$ & & $10.188$ & $10.729$ & & $8.393$ & $8.781$ & & $5.933$ & $6.099$\\
\midrule
fit $a$ & \multicolumn{3}{c|}{$0.0443$ ($b=0.32$, $c=0.68$)} & \multicolumn{3}{c|}{$0.0266$ ($b=0.35$, $c=0.82$)} & \multicolumn{3}{c}{$0.0128$ ($b=0.24$, $c=1.51$)}\\
\bottomrule
\end{tabular}
\end{center}
(three-parameter fits over $p\ge100$; the Rankin--Selberg estimate at $p=1009$ already lies below $\Lambda_1$ for these levels, $X/L<50$, and is omitted beyond). The leading coefficients follow $\frac6{\pi(D+1)}$ across the whole range $D=3,\dots,163$, a factor $41$ between the extremes; but at $D=163$ the term $b\log p$ with $b\approx0.24$ exceeds the leading term for $p<10^9$, and the suppression of the norm is, in the range of any computation, a suppression of $(\log p)^2$ only. Section~\ref{sec:sieve} says what it means for the sieve.

\subsection{The even discriminants $D=4,8$}
For $D=4$ ($\theta_4=\theta^2$, $w_4=4$) and $D=8$ ($\theta_8(z)=\sum e((x^2+2y^2)z)$, $w_8=2$) the theta series are again the Eisenstein series $E_1(z,\chi_{-D})$ of weight one, since $h=1$, and the level $Dp$ has, besides the Atkin--Lehner cusps $W_Q\infty$, $Q\mid Dp$, the cusps of denominator $2\delta$ ($D=4$) and $2\delta,4\delta$ ($D=8$) of Part~II, Section~6.6. The two ingredients that fix the constants are the same: $r_D=w_D\sigma_\chi$ and the class number formula $w_D^2L(1,\chi_{-D})^2=4\pi^2/D$. Hence the Rankin--Selberg constant is $C_D=w_D/(4\pi L(1,\chi_{-D}))=w_D^2\sqrt D/(8\pi^2)$ for all nine discriminants ($C_4=4/\pi^2$, Part~II; $C_8=\sqrt2/\pi^2$, with the same proof once the cusps of $\Gamma_0(8p)$ over the cusp $\infty$ of $\Gamma_0(8)$ are listed, which we have not written out), and the residue of $D_D(s)=\sum r_D(k)^2k^{-s}$ at $s=1$ is $w_D^2L(1,\chi)^2\cdot\frac6{\pi^2}\cdot\prod_{\ell\mid D}(1+\ell^{-1})^{-1}=24/\psi(D)$, where the Euler factor at each ramified prime is $(1+\ell^{-s})^{-1}$ because $\sigma_\chi(\ell^j)=1$ there, and $\psi(D)=D\prod_{\ell\mid D}(1+1/\ell)$ is the Dedekind function. The two wide Atkin--Lehner cusps (widths $Dp$ and $p$, which carry $\pm G_{1,p}$ by the same computation as in Theorem~\ref{thm:AL}, using $\theta_D(-1/(Dz))=-i\sqrt Dz\,\theta_D(z)$ for $D=8$ as well) therefore give, uniformly for the nine discriminants,
\[
\|G^{(D)}_{1,p}\|^2\ \ge\ \frac{6}{\pi\,\psi(D)}(\log p)^2+O(\log p)=\frac{2}{\vol(\Gamma_0(D)\backslash\HH)}(\log p)^2+O(\log p),
\]
since $\vol(\Gamma_0(D)\backslash\HH)=\frac\pi3\psi(D)$; this is $\frac1\pi$ for $D=4$ (Part~II, Theorem~6.14, where it is the asymptotic), $\frac1{2\pi}$ for $D=8$, and $\frac6{\pi(D+1)}$ for the odd $D$. The growth constant of the cusp form is thus the reciprocal of the volume of the modular curve of the \emph{discriminant}, not of the level: the natural normalisation of $\|G\|^2$ against $\vol(\Gamma_0(Dp)\backslash\HH)$, which the Rankin--Selberg mean square $\sum_{k\le Y}|a(k)|^2/k\sim16\pi^2\|G\|^2Y/\vol(\Gamma_0(Dp)\backslash\HH)$ uses, makes the mean square of $r_D(\alpha)r_D(\beta)$ over $\alpha+p\beta\le Y$ equal to $2(\log Y)^2\cdot\frac{Y}{p}\cdot\frac{\vol(\Gamma_0(D)\backslash\HH)}{\vol(\Gamma_0(Dp)\backslash\HH)}\cdots$, and the index $[\Gamma_0(D):\Gamma_0(Dp)]=p+1$ cancels against the $1/p$ of the sparse sequence $p\beta$.

\section{Does a large discriminant help the sieve?}\label{sec:sieve}
For the sieve problem of Parts~I--III with $r_D$ in place of $r$, i.e.\ the sequence $r_D(n)r_D(N-n)$, the norm $\|G^{(D)}_{m,m'}\|^2\sim\frac6{\pi\psi(D)}(\log mm')^2$ enters the remainder through the uniform bound $|a_G(N)|\ll N^{1/2+\e}\|G\|$ (Part~I, Theorem~2.3, whose proof is independent of $D$), and it decreases like $1/D$. This does not help. First, the main term decreases as well: the density of the sequence is governed by $r_D$, whose mean value is $w_DL(1,\chi_{-D})=2\pi/\sqrt D$, so the main term $X_N=\sum_nr_D(n)r_D(N-n)\asymp\frac{4\pi^2}D\,N\,\mathfrak S_D(N)$ is of size $N/D$, while the remainder for one pair $(d,e)$ of sieve moduli is $\ll N^{1/2+\e}(\log de)/\sqrt{\psi(D)}$; the ratio remainder/main is therefore $\asymp\sqrt D\cdot N^{-1/2}(\log de)$ for prime $D$, \emph{larger} by $\sqrt D$, not smaller. Second, and decisively, the level of distribution is an exponent, and $D$ is fixed: the proof of the level $1/2$ (Part~I, Theorem~3.2) sums $|a_{G_{d,e}}(N)|$ over $de\le N^{1/2-\e}$ and loses only constants and $N^\e$ in $D$, so the level is $1/2$ for every $D$, no more and no less; and the sieve dimension is $2$ for every $D$, since $\sum_{p\le z}\omega_D(p)/p=2\log\log z+O(1)$ with $\omega_D(p)=2(1+\chi_{-D}(p))$ on average. What a large discriminant changes is only the constant in the final count ($\asymp N/(D(\log N)^2)$ in place of $N/(\log N)^2$) and, for the odd $D$, the shape of the fundamental domain (no cusps of denominator $2\delta$), which simplifies the analysis of the norm but has no bearing on the dispersion problem of Section~\ref{sec:eis}.

\section{The Eisenstein dispersion form}\label{sec:eis}
Here $D=3$, $m,m'$ coprime square-free integers prime to $3N$, and $T(m,m';N)=\sum_{m\alpha+m'\beta=N}r_3(\alpha)r_3(\beta)$. We follow Part~III, Section~3 (in its corrected form: the Kloosterman sums there are $S(\overline{16m}N,|h|^2;m'k)$, a Kloosterman sum depending on its arguments only through their product, so that the inverse of $16m$ modulo the modulus cannot be removed).

\begin{lemma}[The circle modulo $c$]\label{lem:circle3}
Let $(c,3)=1$, $c$ odd or even, $(a,c)=1$, $h\in\Z^2$, and $|h|_3^2=h_1^2-h_1h_2+h_2^2$ the norm of the dual (hexagonal) lattice. With $c'^{*}$ the product of the primes dividing $c'$ to an odd power and $\chi=\chi_{-3}$ (so $\chi(2)=-1$),
\[
F_c(a;h):=\sum_{\substack{x,y\bmod c\\ N_3(x,y)\equiv a\ (c)}}e\Bigl(\frac{h_1x+h_2y}c\Bigr)=\sum_{g\mid(c,h_1,h_2)}g\,\chi\Bigl(\Bigl(\frac cg\Bigr)^{*}\Bigr)S\Bigl(a,\ \overline3\,\Bigl|\frac hg\Bigr|_3^2;\ \frac cg\Bigr).
\]
\end{lemma}
\begin{proof}
As in Part~III, Lemma~3.1. Detect the congruence by $\frac1c\sum_{t\bmod c}e(t(N_3(x,y)-a)/c)$; for $(t,c)=g$ the sum over $(x,y)$ vanishes unless $g\mid h$, and for $(t,c)=1$ one completes the square: with the Gram matrix $A=\left(\begin{smallmatrix}2&1\\1&2\end{smallmatrix}\right)$, $N_3=\frac12x^{\mathsf T}Ax$, the substitution $x\mapsto x-\bar tA^{-1}h$ (legitimate modulo $c$ since $\det A=3$ is invertible, also at $2$) gives $\sum_{x,y}e((tN_3(x,y)+h\cdot x)/c)=e(-\bar t\,h^{\mathsf T}A^{-1}h/(2c))\,G_A(t;c)$ with $h^{\mathsf T}A^{-1}h/2=|h|_3^2/3$, and $G_A(t;c)=\sum_{x,y}e(tN_3(x,y)/c)=G(t;c)\,G(3\bar4t;c)$ for odd $c$ after diagonalising $N_3=(x+y/2)^2+\frac34y^2$, which equals $c\,\chi((c)^{*})$ by the evaluation of the quadratic Gauss sums prime by prime, the twisting Jacobi symbols of the Chinese remainder theorem squaring to $1$; at the prime $2$ the form is unimodular and $G_A(t;2^\nu)=(-1)^\nu2^\nu=\chi((2^\nu)^{*})2^\nu$. The script \texttt{eisenstein\_circle\_check.py} verifies the identity for all $a$, $h$ and all $c\le40$ coprime to $3$, odd and even ($285\,902$ cases).
\end{proof}

\emph{The decomposition.} Since the ideal $(\sqrt{-3})$ is principal of norm $3$, $r_3(3\beta')=r_3(\beta')$, and the terms of $T(m,m';N)$ with $3\mid\beta$ are $T(m,3m';N)$: with $T^\flat(m,m';N)=\sum_{3\nmid\beta}r_3(\alpha)r_3(\beta)$,
\[
a_{G_{m,m'}}(N)-a_{G_{m,3m'}}(N)=T^\flat(m,m';N)-M^\flat(m,m';N),
\]
the analogue of $a_G(N)-a_G(N/2)=T^{\mathrm{odd}}-M^{\mathrm{odd}}$ of Part~III (there $r(2n)=r(n)$ for the ideal $(1+i)$). In $T^\flat$, $r_3(\beta)\ne0$ forces $\beta\equiv1\pmod3$ (every prime $\equiv2\pmod3$ divides $\beta$ to an even power), so $\chi(\beta/k)=\chi(k)$ and $r_3(\beta)=12\sum_{k\mid\beta,\,k<\sqrt\beta}\chi(k)+6\chi(\sqrt\beta)[\beta=\square]$; the divisors $k$ are prime to $3$ and to $m$ (a common prime would divide $N$) but may be even. Hence
\begin{gather*}
T^\flat(m,m';N)=12\sum_{\substack{k<\sqrt{N/m'}\\ 3\nmid k}}\chi(k)\,\mathcal N_k+6\sum_j\chi(j)\,r_3^\flat\Bigl(\frac{N-m'j^2}m\Bigr),\\
\mathcal N_k=\#\Bigl\{(x,y):\ N_3(x,y)\equiv N\bar m\ (m'k),\ mN_3(x,y)\not\equiv N\ (3),\ 0<N_3(x,y)<Y_k\Bigr\},
\end{gather*}
$Y_k=(N-m'k^2)/m$. The condition modulo $3$ is the analogue of the parity condition $x+y$ odd of Part~III: Poisson summation on the classes modulo $3c$, $c=m'k$, gives
\[
\mathcal N_k=\frac{Y_k}{(3c)^2}\sum_{h\in\Z^2}\hat w\Bigl(\frac{h\sqrt{Y_k}}{3c}\Bigr)\,\epsilon_3(h)\,F_c(N\bar m;\ \bar3h),
\]
where $\bar3h$ is the Chinese-remainder twist and $\epsilon_3(h)$ is the Fourier transform of the admissible classes modulo $3$: since $N_3(x,y)\equiv(x-y)^2\pmod3$, the condition $mN_3\not\equiv N\pmod 3$ is $x\equiv y$ if $N\bar m\equiv1$, $x\not\equiv y$ if $N\bar m\equiv0$, and empty if $N\bar m\equiv2\pmod3$, so that
\[
\epsilon_3(h)=3\,[h_1+h_2\equiv0\ (3)],\qquad 9[h\equiv0\ (3)]-3[h_1+h_2\equiv0\ (3)],\qquad 9,
\]
respectively. As $|h|_3^2\equiv(h_1+h_2)^2\pmod3$, the condition $h_1+h_2\equiv0\pmod3$ is $3\mid|h|_3^2$: exactly the condition $h_1\equiv h_2\pmod2$, i.e.\ $2\mid|h|^2$, of Part~III, with the ramified prime of the field in the role of $2$. By Lemma~\ref{lem:circle3} the leading term $g=1$ of $F_c(N\bar m;\bar3h)$ is $\chi((m'k)^{*})S(N\bar m,\overline{27}|h|_3^2;m'k)$, and $\chi(k)\chi((m'k)^{*})=\chi(m')\chi(k/k^{*})=\chi(m')$ for square-free $m'$, a constant sign for each $m'$ (in Part~III, $m'\equiv1\pmod4$ made it $+1$). We have:

\begin{proposition}\label{prop:eisdisp}
With the notation above, the leading term ($g=1$) of $a_{G_{m,m'}}(N)-a_{G_{m,3m'}}(N)$ is
\[
12\,\chi(m')\sum_{\substack{k<\sqrt{N/m'}\\ 3\nmid k}}\frac{Y_k}{(3m'k)^2}\sum_{h\ne0}\hat w\Bigl(\frac{h\sqrt{Y_k}}{3m'k}\Bigr)\,\epsilon_3(h)\,S\bigl(\overline{27m}\,N,\ |h|_3^2;\ m'k\bigr),
\]
and the remainder sum of the sieve at the level $N^\theta$ has as its leading term the trilinear form $\sum_{m'}\Lambda'_{m'}\sum_{c\equiv0\,(m')}\sum_m\Lambda_m\sum_jb_{m,j}(c)S(\overline{27m}N,j;c)$ over the moduli $c=m'k$ prime to $3$, odd or even, with $j=|h|_3^2\ll mm'$ a norm from $\Z[\zeta_3]$; by the Chinese remainder theorem its $m$-th piece is the bilinear form in the two moduli
\[
\mathfrak B^{(3)}_m=\sum_{m'}\Lambda'_{m'}\sum_{k\le K}\sum_hw_{k,m'}(h)\,S\bigl(\overline{27m}N\bar k,\ |h|_3^2\bar k;\ m'\bigr)\,S\bigl(\overline{27m}N\overline{m'},\ |h|_3^2\overline{m'};\ k\bigr),
\]
the inverse of $27m$ taken modulo $m'$ and modulo $k$ respectively.
\end{proposition}

\emph{Comparison with the Gaussian case.} Nothing in the structure changes. (i) The fixed argument is again the single integer $N$, carried by a unit ($\overline{27m}$ in place of $\overline{16m}$) modulo the modulus, so the obstacle (a) of Part~III, Section~3.4, the absence of an average over it, is the same, and so is the obstruction to the Kuznetsov formula. (ii) The sizes are the same: $|h|\ll m'k\sqrt{m/N}$, so the $h$-box has side $H\le\sqrt{mm'}$, below $\sqrt q$ for every modulus $q=[c_1,c_2]$ of a Poisson step, and the Cauchy--Schwarz arrangements (A1), (A2) of Part~III, Section~4 give the same diagonals ($E^{1/4}N^{3/4}$ and $\sqrt{EN}$); there is no tighter grouping, because the groupings depend only on the lengths of the variables and on the Weil mean square of Kloosterman sums over the dual lattice, which is the same for the hexagonal lattice. (iii) The modulus $m'k$ now ranges over all $k$ prime to $3$, even $k$ included: more moduli, no gain. (iv) The $2$-adic ``impurity'' of Part~III (the cusps of denominator $2\delta$, the parity condition) has moved to the prime $3$: the condition $3\nmid\beta$ produces $\epsilon_3(h)$ and the dual condition $3\mid|h|_3^2$, as $2\nmid\beta$ produced $h_1\equiv h_2\pmod2$. This is not an accident: the ramified prime of the field is where the ideal of norm $\ell$ is principal, and removing it from the divisor sum of $r_D(\beta)$ is what makes the divisors $k<\sqrt\beta$ pair with $\beta/k$ under $\chi$. The clean level $3p$ matters for the holomorphic side (Sections~2--5), not for the dispersion form; a level beyond $1/2$ needs the bilinear saving over $(m',k)$ of Part~III, Section~4.5, in the Eisenstein setting as in the Gaussian one.

\section*{Scripts}
\texttt{code/eisenstein\_circle\_check.py} (Lemma~\ref{lem:circle3}, and the Kloosterman product identity behind the correction of Part~III; output in \texttt{results/eisenstein\_circle\_check.txt}).

\section*{Script}
\texttt{code/heegner\_theta\_unified.py} (self-contained; Kronecker symbol, $r_D$, $T$, $M$, counts of local densities, Atkin--Lehner check at a base point, Hecke decomposition by $T_2$ then $T_3,T_5,T_7,T_{13}$ where eigenvalues coincide, central values by the smoothed functional equation, norms by Rankin--Selberg and by the cusp sums); output in \texttt{results/heegner\_theta\_unified.txt}. Floating point.
\end{document}
